% =============================================================
% §7 相关工作
% =============================================================
\section{相关工作}
\label{sec:related}

\textbf{LLM 服务系统（单机或集群内）。}Orca~\cite{yu2022orca} 引入迭代级连续批处理；vLLM~\cite{kwon2023vllm} 加入 PagedAttention 以实现非连续 KV 分配（吞吐 $2$--$4\times$）；Sarathi-Serve~\cite{agrawal2024sarathi} 引入分块 prefill 以缓解混合批中的生成阻塞。VoltanaLLM~\cite{yu2025voltana} 在 prefill/decode 解耦的集群内部，通过频率控制与状态空间路由协同优化能耗与 \slo{}。这些系统优化的是单个服务引擎或集群\emph{内部}的行为，与本文的跨服务器调度\textbf{正交}：本文的准入式批处理是它们迭代级批处理的一个可计算抽象，本文的逐请求放置决策可无改地传入 vLLM 风格引擎。

\textbf{集群调度：从启发式到学习。}经典调度器——按上行 rank 排定优先级的 HEFT~\cite{topcuoglu2002heft}，以及 GA~\cite{holland1975adaptation}、PSO~\cite{kennedy1995pso} 等元启发式——针对通用 DAG 负载优化固定代价模型；产业界则偏好更简单的负载均衡原语。学习式调度器随之出现：Decima~\cite{mao2019decima}（Spark DAG 上的 GNN + REINFORCE）、RLScheduler~\cite{zhang2020rlscheduler}（HPC 批作业），以及 A3C-R2N2~\cite{tuli2022a3cr2n}（面向通用边缘--云任务的 actor-critic，早于 LLM 时代）——后者是部署拓扑上最接近的前驱，我们将其作为\textbf{对照基线}进行基准比较（\S\ref{sec:main-results}）以隔离本文的 AIGC 贡献。这一路线至今活跃：近期工作将基于 GNN 的 DRL 用于能耗感知云调度~\cite{dagtopology2026}，将扩散辅助的奖励塑形用于跨数据中心的 AIGC 负载调度~\cite{aigcenergy2026}。该文献中的两个发现预示了本文结论：Decima 的收益集中在高集群负载下，而 \cite{dagtopology2026} 观察到每个学得策略都会退化为某一 regime 特定的启发式——与本文的 regime 刻画（C4）一致。这些调度器无一编码 LLM 服务物理。

\textbf{跨服务器的 LLM 请求调度。}一条快速增长的路线在实例\emph{之间}而非单引擎之内调度 LLM 推理。在单集群内部，Microsoft 的负载感知路由器~\cite{jain2024router} 在 prefill/decode 与批混合效应上训练启发式引导的 RL 策略；Lodestar~\cite{lim2026lodestar} 学习一个将批处理与 KV cache 复用耦合的在线奖励预测器——两者都呼应了本文"应把服务物理暴露给学习式路由器"的论点，但均未建模云--边层级或能耗。FREESH~\cite{he2025freesh} 通过经典优化在地理分布的异构 GPU 上为能耗与碳排放做路由与调度；MARLIN~\cite{moore2026marlin} 将多智能体博弈论 RL 应用于跨地理分布云数据中心的多目标 LLM 调度（无边缘层）。在云--边一侧，PerLLM~\cite{yang2024perllm} 是最接近的已有工作：它用带约束的组合 bandit 在边缘与云之间调度多样化的 LLM 服务请求，并报告了可观的吞吐与能耗成本收益。本文在三个轴上与之不同：\emph{（i）方法}——PerLLM 的 bandit 在约束下优化一个标量化代价，而本文学习一个\textbf{状态与奖励显式暴露服务物理}（模型驻留、KV cache 本地性、批占用）的深度策略；\emph{（ii）目标}——本文将 \slo{} 与每 token 能耗视作显式的 Pareto 目标对，而非折叠为单一代价；\emph{（iii）分析}——本文的联合消融与权重敏感性研究（\S\ref{sec:ablation}、\S\ref{sec:weight-sens}）量化了此类增益\emph{从何而来}，这是 bandit 类方法无法给出的分析。我们在仿真器中重新实现了一个 CS-UCB 风格的调度器作为基线；它在 Pareto 平面上与通用基线聚为一簇（\S\ref{sec:main-results}）。其他云--边工作在决策粒度或方法上有所不同：Lyapunov 辅助的 DRL 在层级之间以子层粒度切分 transformer 计算~\cite{younesi2025lyapunov}；主动推断卸载~\cite{he2024activeinference} 与带端侧压缩的世界模型辅助 PPO 卸载~\cite{zhang2026worldmodel} 做设备侧卸载决策；T2DRL~\cite{xu2025t2drl} 在移动边缘协同设计模型缓存与卸载；QoS 感知路由在边缘部署的专家之间做选择~\cite{yang2025qosrouting}；NSGA-II 搜索在云--边 LLM 实例之间路由请求而不含学习~\cite{yu2025llmrouting}——我们将其实现为本文最强基线（\S\ref{sec:eval}）；SLICE~\cite{chow2025slice} 则在单台边缘设备上通过效用启发式与解码速率控制调度差异化的 \slo{} 层级。据我们所知，\textbf{尚无已有工作将异构云--边池上的深度 RL 策略、服务物理状态/奖励与 \slo{}+能耗联合 Pareto 评估结合起来}——这正是本文所研究并消融的具体组合。

\textbf{AIGC 部署拓扑。}Splitwise~\cite{patel2024splitwise} 与 DistServe~\cite{zhong2024distserve} 将 prefill 与 decode 解耦到不同代次的机器上；Llumnix~\cite{sun2024llumnix} 在实例之间中途迁移请求。这些工作设计的是\emph{部署拓扑}与\emph{反应式策略}，将初始请求放置视为正交问题——两条线索可以组合：本文调度器可为 Splitwise 风格部署填充 prefill/decode 池，而 Llumnix 风格的迁移可在运行时精化本文的放置。

\textbf{仿真平台。}CloudSim~\cite{calheiros2011cloudsim}、iFogSim~\cite{gupta2017ifogsim} 与 WorkflowSim~\cite{chen2012workflowsim} 针对通用云--雾或 DAG 负载，不含 LLM 物理。近期涌现的一批 LLM 服务仿真器以高保真度建模集群内服务：Vidur~\cite{agrawal2024vidur} 构建"性能剖析 + 机器学习"的算子模型以支持集群规模的 what-if 分析；SplitwiseSim~\cite{patel2024splitwise} 仿真阶段拆分集群；LLMServingSim~\cite{cho2024llmservingsim} 软硬件协同仿真；Frontier~\cite{frontier2025sim} 以显式的跨服务器 KV 传输事件建模解耦服务；Helix 的仿真器覆盖地理分布的异构 GPU 池~\cite{mei2025helix}。Vidur 的一个扩展耦合了 GPU 功耗模型以量化能耗与碳排放~\cite{wiesner2025vidurvessim}。因此，本文仿真器的每一根支柱都各自孤立地存在于某处——连续批处理与两阶段生命周期已是标准配置；Frontier 建模跨服务器 KV 移动；\cite{wiesner2025vidurvessim} 加入能耗核算；权重驻留/冷加载动力学则出现在 serverless 服务系统中~\cite{fu2024serverlessllm}。然而据我们所知，本文平台是首个将全部五项 AIGC 物理、云--边跨服务器调度范围与原生多目标（\slo{}~+~能耗）评估\textbf{结合}起来的\emph{开源}平台。我们随 \S\ref{sec:eval} 的 30+ 配置一并发布可复现 manifest。
